\completa{12}{Sergio}{5a.1: por qué agrupamos antes de filtrar}
\senalar{Seguir los cuatro cuadros y señalar la posición de filter después de agg.}
\decir{El punto cinco a pide agrupar y filtrar utilizando dos o más columnas. Departamento y rango de edad nos permiten preguntar: ¿qué combinaciones reúnen al menos mil personas reportadas? Buscamos acumulados por perfil territorial, no filas individuales con una cantidad alta. El umbral de mil es una condición de selección de esta consulta, no un criterio oficial de riesgo.}
\decir{El primer cuadro dice «noventa y tres mil ochocientos treinta y un registros». La flecha conduce a «groupBy: departamento y edad». GroupBy forma un grupo por cada pareja de valores. Por ejemplo, Lima Metropolitana y doce a diecisiete años es un grupo distinto de Lima Metropolitana y dieciocho a veintinueve. La siguiente flecha lleva a «SUM de cantidad; filter mayor o igual a mil». Sum suma cantidades dentro de cada grupo; filter conserva los grupos que cumplen el umbral. La última flecha llega a «treinta combinaciones; orden descendente»: primero aparecen los totales mayores.}
\decir{El código expresa esa misma secuencia. Denuncias es el DataFrame limpio. Agg aplica la agregación; as personas da un nombre a la suma. Col personas referencia esa nueva columna. OrderBy con desc ordena de mayor a menor. Val consulta guarda la definición del resultado en una variable de Scala; los puntos encadenan operaciones y los paréntesis contienen sus argumentos.}
\decir{En la tabla, Lima Metropolitana de doce a diecisiete años suma diecisiete mil novecientas una personas; de dieciocho a veintinueve, siete mil ochocientas setenta y cinco; y Junín de doce a diecisiete, cuatro mil ciento sesenta y nueve. Es una muestra de las treinta combinaciones devueltas. El filtro se aplica al acumulado ya calculado. Una fila con cantidad uno puede aportar a un grupo que supere mil. Ahora mantendremos ese principio y cambiaremos la pregunta hacia los distritos.}

\completa{13}{Sergio}{5b.2: del acumulado al ranking}
\senalar{Mostrar que orderBy precede a limit y que la clave territorial tiene tres componentes.}
\decir{El punto cinco b solicita un orden decreciente. La disponibilidad de departamento, provincia y distrito permite preguntar: ¿cuáles son los quince distritos con más personas reportadas entre 2019 y 2024? Las tres columnas evitan mezclar distritos que comparten nombre en ubicaciones distintas.}
\decir{El primer cuadro dice «Todos los registros». Su flecha llega a «groupBy: departamento, provincia, distrito». La siguiente conduce a «SUM de cantidad; orderBy descendente». Primero calculamos el acumulado por ubicación y luego lo ordenamos de mayor a menor. La última flecha llega a «limit quince; quince distritos». Limit restringe la salida a los primeros quince resultados; debe ir después del orden para que sean los de mayor acumulado.}
\decir{El fragmento Scala utiliza agg para sumar, as para nombrar personas, desc para indicar orden descendente y limit para seleccionar los quince. En la salida, San Juan de Lurigancho suma cuatro mil setecientas setenta y cinco personas; Lima, tres mil ciento cincuenta y tres; San Martín de Porres, dos mil seiscientas cuarenta y seis; y Callao, dos mil seiscientas cuarenta y dos. Ordenamos personas contabilizadas, no número de filas. El ranking describe volumen reportado y no riesgo por habitante. Para agrupar territorios de otra forma necesitamos unirlos con un catálogo; ahora explicaré esa unión.}

\completa{14}{Sergio}{5c.1: INNER JOIN y conservación del total}
\senalar{Seguir las cuatro cajas y detenerse en la clave ubigeo antes de leer las regiones.}
\decir{El punto cinco c pide distintos JOIN. En este primero preguntamos: ¿cómo se distribuye el total entre las regiones definidas por el catálogo geográfico del proyecto? La tabla de denuncias aporta las cantidades y el catálogo aporta la región asociada al ubigeo. Un JOIN relaciona ambas tablas mediante una clave común.}
\decir{El primer cuadro dice «groupBy ubigeo; suma de personas». Reducimos las filas a un acumulado por código geográfico y conservamos también el conteo de registros. La flecha conduce a «INNER JOIN; clave ubigeo». Inner significa que se conservan las claves que encuentran coincidencia en ambos lados. En el código, join recibe dimGeografia, Seq ubigeo y la palabra inner. Seq indica aquí la lista de columnas utilizadas como clave.}
\decir{La siguiente flecha llega a «groupBy región; sumar personas y registros». Una vez incorporada la región, reunimos los ubigeos de cada región. La última flecha conduce a «tres regiones; total ciento veintiún mil cuatrocientas veintisiete». La tabla muestra Costa con setenta y nueve mil seiscientas sesenta y seis personas y quinientos sesenta y siete ubigeos; Sierra con treinta y cuatro mil trescientas una y ochocientos cuatro; Selva con siete mil cuatrocientas sesenta y ciento treinta y cinco. Ubigeos es número de claves geográficas observadas, no número de personas.}
\decir{Para que el resultado sea válido, el catálogo tiene una fila por ubigeo. Si una clave encontrara dos filas de dimensión, su cantidad se repetiría al unir. Además, la salida guardada no presenta claves nulas, es decir, ausentes, en el control geográfico. Verificamos la suma final para comprobar que no perdimos ni multiplicamos el total. La clasificación regional es la del código por departamento; no mide ruralidad. Continuaremos con funciones estadísticas sobre esos registros.}

\completa{15}{Sergio}{5d.3: promedio, desviación y varianza}
\senalar{Leer los cuatro cuadros y las funciones de los dos fragmentos. Distinguir desviación poblacional y varianza muestral.}
\decir{El promedio global de tres b puede ocultar diferencias entre categorías. Por eso, usando SEXO y CANTIDAD, preguntamos: ¿cómo varía la dispersión de CANTIDAD por registro entre mujeres y hombres? Dispersión significa cuánto se separan los valores alrededor de su promedio. Esta consulta ejemplifica el punto cinco d, funciones de Spark, y conserva también NO ESPECIFICA.}
\decir{El primer cuadro dice «Registros por sexo». La flecha lleva a «groupBy sexo»: formamos los tres grupos. La segunda flecha conduce a «avg, stddev pop, variance, min, max». Avg obtiene la media; stddev pop obtiene la desviación estándar poblacional; variance calcula la varianza muestral; min y max encuentran los extremos. El código añade count para filas, sum para personas y round para redondear, en este caso a cuatro decimales. As asigna nombres a las columnas calculadas.}
\decir{La última flecha llega a «tres grupos; medias y dispersiones». En la salida F presenta promedio uno coma tres siete uno cinco y desviación poblacional uno coma cero cuatro cinco cuatro. M presenta uno coma uno ocho seis seis y cero coma cinco nueve siete uno. NO ESPECIFICA presenta uno coma cero seis siete seis y cero coma tres cero cero cero. Una desviación menor indica conteos más cercanos entre sí, en la escala de CANTIDAD.}
\decir{Conviene distinguir los denominadores: la desviación poblacional se obtiene a partir de una varianza dividida entre n, el número de filas del grupo. La varianza muestral divide entre n menos uno. Además, la varianza está en unidades al cuadrado y la desviación vuelve a la escala original mediante la raíz. Por eso no debemos comparar directamente variance con stddev pop ni elevar al cuadrado la cifra mostrada esperando igualdad exacta. La mayor dispersión de F describe conteos agregados; no demuestra desapariciones en grupo. Ahora explicaré cómo esas mismas tablas pueden consultarse mediante vistas SQL.}

\completa{16}{Sergio}{5e.2: dos JOIN, perfil y porcentaje nacional}
\senalar{Leer las cuatro cajas; marcar ambos JOIN y después el OVER vacío del porcentaje.}
\decir{Al conocer los rangos de edad y el sexo surge la pregunta: ¿qué combinación de grupo etario y sexo concentra más personas y qué porcentaje del total representa? Grupo etario significa categoría de edad, como adolescente. Esta consulta resuelve cinco e mediante vistas temporales y JOIN. Usamos las vistas denuncias, dim edad y dim sexo; las claves son rango edad y sexo.}
\decir{El primer cuadro dice «Vistas registradas en esquema punto Scala». Registrar significa dar a cada DataFrame un nombre que spark punto sql puede consultar. CreateOrReplaceTempView crea esa vista o sustituye su definición en la sesión; no genera una tabla permanente de PostgreSQL. La flecha conduce a «JOIN dim edad; JOIN dim sexo». La primera unión incorpora el nombre del grupo etario mediante rango edad. La segunda relaciona la categoría de sexo. Son tres tablas y dos operaciones JOIN.}
\decir{La siguiente flecha llega a «GROUP BY grupo y sexo; suma y porcentaje nacional». Group By forma las combinaciones y SUM calcula su cantidad. El código utiliza d, e y s como alias, nombres cortos de denuncias, edad y sexo; por eso d punto cantidad indica de qué tabla viene la medida. ON expresa la condición de igualdad entre claves. SELECT elige las columnas de salida y AS les asigna nombre.}
\decir{El porcentaje divide cada suma entre SUM de SUM de cantidad, OVER paréntesis vacíos. La suma interior corresponde a cada grupo; la suma exterior es una función de ventana que reúne esos totales sin eliminar las filas del resultado. OVER vacío toma todos los grupos, por eso el denominador es nacional. CAST AS DOUBLE convierte a un tipo numérico con decimales y ROUND redondea a dos posiciones. El código incluye también COUNT asterisco para contar registros y AVG para obtener el promedio de CANTIDAD por registro.}
\decir{La última flecha lleva a «dieciocho combinaciones; orden por personas». Adolescentes mujeres reúne cincuenta y un mil veintiuna personas, cuarenta y dos coma cero dos por ciento del total nacional. Adolescentes hombres reúne trece mil trescientas nueve, diez coma noventa y seis por ciento. Adultos hombres reúne doce mil trescientas ochenta y nueve, diez coma veinte por ciento. La fórmula visible multiplica cincuenta y un mil veintiuna por cien y divide entre ciento veintiún mil cuatrocientas veintisiete. Es participación dentro del dataset, no probabilidad individual. Usaremos ahora filtros para delimitar otro universo de comparación.}

\completa{17}{Sergio}{5f.1: el orden de WHERE, GROUP BY, HAVING y LIMIT}
\senalar{Seguir los cuatro cuadros; señalar dónde se filtran filas y dónde se filtran grupos.}
\decir{Después del ranking general podemos hacer una pregunta más específica: ¿qué diez distritos reúnen más mujeres reportadas en 2024 con rango de edad especificado? El punto cinco f pide ORDER BY combinado con filtros. Aquí fijamos año, sexo y calidad de la información de edad, y exigimos al menos veinte personas por distrito.}
\decir{El primer cuadro dice «WHERE: año, sexo y edad». Where selecciona las filas que cumplen año igual a 2024, sexo igual a F y rango de edad distinto de NO ESPECIFICA. El signo menor que seguido de mayor que en SQL significa distinto. La primera flecha llega a «GROUP BY: departamento, provincia, distrito». Con las filas ya filtradas formamos cada ubicación y sumamos CANTIDAD.}
\decir{La siguiente flecha conduce a «HAVING SUM mayor o igual a veinte; ORDER BY personas descendente». Having filtra grupos ya sumados; Order By los ordena de mayor a menor. La última flecha lleva a «LIMIT diez; diez distritos». Seleccionamos las diez primeras ubicaciones del ranking después de filtrar y ordenar. En el código, COUNT asterisco añade el número de registros, mientras SUM aporta personas.}
\decir{La salida muestra San Juan de Lurigancho con trescientas setenta y nueve; Lima con doscientas ochenta y cuatro; Ate con doscientas setenta y una; y Callao con doscientas sesenta y cinco. No deben compararse directamente con los cuatro mil setecientos setenta y cinco del ranking general: este universo comprende solo mujeres de 2024 con edad especificada. Tampoco sería equivalente aplicar cantidad mayor o igual a veinte en WHERE: eso buscaría filas individuales de ese tamaño y no distritos con veinte personas acumuladas. Ya tenemos las respuestas analíticas y la base descrita; ahora mostraré cómo se entregan estos resultados en Excel.}
